\kefalaio{Χαρακτήρες διαιρετότητας: ξέρω χωρίς να διαιρέσω}{Κεφάλαιο Α.1.5, «Χαρακτήρες διαιρετότητας – ΜΚΔ – ΕΚΠ – Ανάλυση αριθμού σε γινόμενο πρώτων παραγόντων» (πρώτο μέρος)}

\begin{anagke}
Το σχολικό παζάρι πήγε καλύτερα απ' όσο περίμενε κανείς. Η Φανή, που κρατάει το ταμείο, μετράει τα χρήματα: 2.358 €. Θέλουν να τα μοιράσουν ίσα σε τρεις ομάδες που θα φτιάξουν από ένα σχολικό έργο. Μετά η Μυρτώ λέει: «Και αν οι ομάδες ήταν εννιά;» Ο Άρης: «Και αν ήταν πέντε;» Ο Στέλιος ξεκινάει να κάνει τις διαιρέσεις με τη σειρά, μία μία. Τρεις διαιρέσεις με τετραψήφιο αριθμό είναι πολλή δουλειά.

Η Λίνα κοιτάζει τον αριθμό 2.358 και λέει: «Για το πέντε δεν χρειάζεται καν να κάνεις διαίρεση. Τελειώνει σε οκτώ. Δεν γίνεται». «Πώς το ξέρεις;» ρωτάει ο Ηλίας. «Το έχω δει σε όλους τους αριθμούς που διαιρούνται με το πέντε: τελειώνουν σε 0 ή σε 5». Η Ζωή ρωτάει: «Και για το τρία και για το εννιά υπάρχει κάτι τέτοιο;»

Υπάρχει. Και είναι από τα πιο χρήσιμα κόλπα στο γυμνάσιο: να καταλαβαίνεις αν ένας αριθμός διαιρείται με έναν άλλον κοιτάζοντας απλώς τα ψηφία του.
\end{anagke}

\begin{lexi}{κριτήριο, χαρακτήρας διαιρετότητας}{κριτήριο, χαρακτήρας}
Το \lexiword{κριτήριο} έρχεται από το «κρίνω», δηλαδή αποφασίζω. Κριτήριο είναι ο κανόνας με τον οποίο κρίνουμε κάτι. Ένας κριτής σε αγώνα έχει κριτήρια για να βγάλει τον νικητή. Ένα \lexiword{κριτήριο διαιρετότητας} μας βοηθάει να κρίνουμε αν ένας αριθμός διαιρείται με κάποιον άλλον, χωρίς να κάνουμε τη διαίρεση.

Το σχολικό βιβλίο τα λέει και \lexiword{χαρακτήρες διαιρετότητας}. Ο χαρακτήρας είναι το «σημάδι» που χαρακτηρίζει κάτι, όπως ο χαρακτήρας ενός ανθρώπου. Εδώ το σημάδι είναι στα ψηφία του αριθμού.
\end{lexi}

\begin{metafrasi}
Ο αριθμός είναι άρτιος. & Λήγει σε $0,\,2,4,\,6$ ή $8$ (διαιρείται με το 2)\\
Ο αριθμός διαιρείται με το 5. & Λήγει σε $0$ ή $5$\\
Ο αριθμός διαιρείται με το 10. & Λήγει σε $0$\\
Ο αριθμός διαιρείται με το 3. & Το άθροισμα των ψηφίων του διαιρείται με το 3\\
Ο αριθμός διαιρείται με το 9. & Το άθροισμα των ψηφίων του διαιρείται με το 9\\
Ο αριθμός διαιρείται με το 4. & Ο αριθμός από τα δύο τελευταία ψηφία διαιρείται με το 4\\
Ο αριθμός διαιρείται με το 25. & Λήγει σε $00$, $25$, $50$ ή $75$\\
Ο αριθμός διαιρείται με το 6. & Διαιρείται και με το 2 και με το 3\\
\end{metafrasi}

\tmima{Η θεωρία}

\begin{kanonas}{κριτήρια διαιρετότητας}
\begin{itemize}[leftmargin=1.2em,itemsep=0.3ex,topsep=2pt]
\item \textbf{Με το 2}: το τελευταίο ψηφίο είναι 0, 2, 4, 6 ή 8.
\item \textbf{Με το 5}: το τελευταίο ψηφίο είναι 0 ή 5.
\item \textbf{Με το 10}: το τελευταίο ψηφίο είναι 0. (Με το 100: τα δύο τελευταία ψηφία είναι 00, με το 1.000: τα τρία τελευταία, κ.ο.κ.)
\item \textbf{Με το 4}: ο αριθμός που σχηματίζουν τα δύο τελευταία ψηφία διαιρείται με το 4 (ή είναι 00).
\item \textbf{Με το 25}: τα δύο τελευταία ψηφία είναι 00, 25, 50 ή 75.
\item \textbf{Με το 3}: το άθροισμα των ψηφίων διαιρείται με το 3.
\item \textbf{Με το 9}: το άθροισμα των ψηφίων διαιρείται με το 9.
\end{itemize}
\end{kanonas}

Από τα κριτήρια αυτά φτιάχνουμε και άλλα. Για να διαιρείται ένας αριθμός με το 6 πρέπει να διαιρείται και με το 2 και με το 3. Για το 15, και με το 3 και με το 5. Για το 12, και με το 3 και με το 4.

\tmimasmall{Γιατί δουλεύουν}

Κάθε αριθμός ξεκινάει με πολλαπλάσια του 10, του 100, του 1.000. Για παράδειγμα $3.725=37\cdot 100+25$. Ο 100 διαιρείται με το 4 και με το 25, άρα το $37\cdot 100$ διαιρείται και με τους δύο, και το μόνο που μένει να ελέγξουμε είναι τα δύο τελευταία ψηφία, το 25. Το ίδιο και για τα 2 και 5, με το 10.

Για το 3 και το 9 δουλεύει ένα κόλπο. Οι αριθμοί 9, 99, 999, \dots\ διαιρούνται με το 9 (και με το 3). Γράφω $2.457=2\cdot 1.000+4\cdot 100+5\cdot 10+7$ και αντικαθιστώ το $1.000=999+1$, $100=99+1$, $10=9+1$:
$$2.457=2\cdot 999+4\cdot 99+5\cdot 9+(2+4+5+7).$$
Τα τρία πρώτα κομμάτια διαιρούνται με το 9. Άρα όλος ο αριθμός διαιρείται με το 9 ακριβώς όταν διαιρείται το $2+4+5+7=18$. Και διαιρείται.

\begin{prosoxi}
\begin{itemize}[leftmargin=1.2em,itemsep=0.3ex,topsep=2pt]
\item Αν ένας αριθμός τελειώνει σε 4, δεν σημαίνει ότι διαιρείται με το 4. Το 14 δεν διαιρείται. Για το 4 κοιτάμε τα δύο τελευταία ψηφία μαζί.
\item Για το 3 και το 9 μετράνε \textbf{όλα} τα ψηφία, όχι μόνο το τελευταίο.
\item Για να διαιρείται ένας αριθμός με το 24 δεν φτάνει να διαιρείται με το 4 και με το 6. Το 12 διαιρείται και με τα δύο, αλλά όχι με το 24. Μπορούμε να «σπάσουμε» έναν αριθμό σε δύο μόνο όταν οι δύο αριθμοί δεν έχουν κοινό παράγοντα (όπως το 3 και το 4, ή το 2 και το 5).
\item Το 0 διαιρείται με κάθε αριθμό (εκτός από το 0 ως διαιρέτη), γιατί $0=\delta\cdot 0$.
\end{itemize}
\end{prosoxi}

\tmima{Λυμένα παραδείγματα}

\paradeigma{Ο πίνακας των κριτηρίων}
\tixero Οι αριθμοί 12.510, 772, 225 και 13.600.\par
\tithelo Με ποιους από τους 2, 3, 4, 5, 9, 10, 25 και 100 διαιρείται ο καθένας.\par
\possindeo Εφαρμόζω κάθε κριτήριο. Στα 3 και 9 υπολογίζω το άθροισμα ψηφίων. Στα 4 και 25 κοιτάζω τα δύο τελευταία ψηφία.

\begin{center}
\small\renewcommand{\arraystretch}{1.3}\setlength{\tabcolsep}{4pt}
\begin{tabular}{@{}lcccccccc@{}}
\toprule
 & 2 & 3 & 4 & 5 & 9 & 10 & 25 & 100\\
\midrule
12.510 \ ($1+2+5+1+0=9$) & ναι & ναι & όχι & ναι & ναι & ναι & όχι & όχι\\
772 \ ($7+7+2=16$, και $72:4=18$) & ναι & όχι & ναι & όχι & όχι & όχι & όχι & όχι\\
225 \ ($2+2+5=9$) & όχι & ναι & όχι & ναι & ναι & όχι & ναι & όχι\\
13.600 \ ($1+3+6=10$) & ναι & όχι & ναι & ναι & όχι & ναι & ναι & ναι\\
\bottomrule
\end{tabular}
\end{center}

\apantisi Δες τον πίνακα.\par
\elegxos Το 13.600 διαιρείται με το 100, άρα και με το 10, το 4, το 25, το 5 και το 2, όπως πρέπει: κάθε διαιρέτης του 100 μπαίνει αυτόματα. Δοκιμή με το 4: $13.600:4=3.400$. Σωστά.

\paradeigma{Το ψηφίο που λείπει}
\tixero Ο αριθμός $4.7\square 2$, όπου το $\square$ είναι ένα ψηφίο.\par
\tithelo Ποιο ψηφίο πρέπει να μπει ώστε ο αριθμός να διαιρείται με το 9.\par
\possindeo Το άθροισμα των γνωστών ψηφίων είναι $4+7+2=13$. Πρέπει $13+\square$ να διαιρείται με το 9. Τα πολλαπλάσια του 9 κοντά στο 13 είναι το 18, άρα $\square=5$. (Το επόμενο πολλαπλάσιο, 27, θα απαιτούσε ψηφίο 14, που δεν υπάρχει.)\par
\apantisi $4.752$.\par
\elegxos $4.752:9=528$, γιατί $9\cdot 528=4.752$. Σωστά.

\paradeigma{Η μοιρασιά του παζαριού}
\tixero Η Φανή έχει 2.358 €.\par
\tithelo Σε ποιες από τις 3, 9, 5, 4 και 6 ισόποσες ομάδες μοιράζονται χωρίς να περισσεύει τίποτα.\par
\possindeo Άθροισμα ψηφίων: $2+3+5+8=18$. Διαιρείται με το 3 και με το 9. Τελευταίο ψηφίο 8, άρα όχι με το 5. Τελευταία δύο ψηφία 58: το $58:4$ δεν είναι ακέραιο, άρα όχι με το 4. Είναι άρτιος και διαιρείται με το 3, άρα διαιρείται με το 6.\par
\apantisi Σε 3, 6 και 9 ομάδες.\par
\elegxos $2.358:3=786$, $2.358:6=393$, $2.358:9=262$. Και οι τρεις είναι τέλειες διαιρέσεις.

\tmima{Ασκήσεις}

\zoni{Το κατάλαβα}{εφαρμογή του κανόνα}

\ex Ποιοι από τους αριθμούς 3.140, 2.025, 1.234, 5.050 και 707 διαιρούνται με το 2; Ποιοι με το 5; Ποιοι με το 10;

\ex Βρες ποιοι από τους 123, 4.521, 1.001, 99.999 και 2.468 διαιρούνται με το 3. Έλεγξε μία από τις περιπτώσεις με διαίρεση.

\ex Βρες ποιοι από τους 1.089, 4.527, 81.081, 7.777 και 90.009 διαιρούνται με το 9.

\ex Βρες ποιοι από τους 316, 1.002, 7.500, 2.024 και 98 διαιρούνται με το 4.

\ex Βρες ποιοι από τους 175, 2.350, 4.600, 1.065 και 3.025 διαιρούνται με το 25.

\ex Συμπλήρωσε έναν πίνακα με τους αριθμούς 840, 1.575, 6.048 και 20.250 και τους διαιρέτες 2, 3, 4, 5, 9, 10, 25. Για κάθε αριθμό γράψε ναι ή όχι.

\ex Ποιοι από τους 54, 128, 4.230, 1.005 και 7.776 διαιρούνται με το 6;

\ex Βρες το ψηφίο (ή τα ψηφία) που πρέπει να μπουν στη θέση του $\square$: α) $4.7\square 2$ να διαιρείται με το 9 \quad β) $3\square 5$ να διαιρείται με το 3 \quad γ) $6\square 4$ να διαιρείται με το 4 \quad δ) $17\square$ να διαιρείται και με το 2 και με το 5.

\ex Σωστό ή λάθος; α) Κάθε αριθμός που τελειώνει σε 4 διαιρείται με το 4. \quad β) Αν ένας αριθμός διαιρείται με το 9, διαιρείται και με το 3. \quad γ) Το 0 διαιρείται με το 7. \quad δ) Αν ένας αριθμός διαιρείται με το 4 και με το 6, διαιρείται και με το 24. \quad ε) Αν διαιρείται με το 2 και με το 5, διαιρείται με το 10. \quad στ) Αν διαιρείται με το 3 και με το 4, διαιρείται με το 12.

\zoni{Το χρησιμοποιώ}{προβλήματα από τη ζωή}

\ex \textbf{Τα χρήματα του παζαριού.} Δες το παράδειγμα της Φανής, αλλά τώρα τα χρήματα είναι 3.465 €. Σε ποιες από τις 3, 5, 9, 4 και 10 ομάδες μοιράζονται ακριβώς; Βρες πόσα παίρνει η κάθε ομάδα, όπου γίνεται.

\ex \textbf{Τα μπλουζάκια.} Ο προπονητής του Άρη έχει 140 μπλουζάκια και θέλει να τα βάλει σε ίσες στοίβες. Με ποιους από τους αριθμούς 2, 3, 4, 5, 6 και 10 γίνεται, χωρίς να περισσεύει κανένα; Χρησιμοποίησε κριτήρια.

\ex \textbf{Ο κωδικός της Λίνας.} Ο νέος κωδικός της Λίνας είναι τετραψήφιος. Ξεκινάει με 7, το δεύτερο ψηφίο είναι το 2, τελειώνει σε 0, διαιρείται με το 9 και το τρίτο ψηφίο είναι μεγαλύτερο από το δεύτερο. Ποιος είναι ο κωδικός;

\ex \textbf{Οι θέσεις του λεωφορείου.} Οι θέσεις ενός λεωφορείου είναι αριθμημένες από το 1 ως το 60. Οι θέσεις που διαιρούνται με το 5 είναι δίπλα στο παράθυρο. Πόσες είναι; Πόσες διαιρούνται με το 10; Πόσες διαιρούνται και με το 3 και με το 5;

\ex \textbf{Τα δίσεκτα έτη.} Ένα έτος είναι δίσεκτο αν διαιρείται με το 4. Εξαίρεση οι αιώνες (1.900, 2.000, \dots), που είναι δίσεκτοι μόνο όταν διαιρούνται με το 400. Ποια από τα έτη 1.900, 2.000, 2.026, 2.028 και 2.100 είναι δίσεκτα;

\zoni{Το σκέφτομαι}{μια ερώτηση που πάει πιο μακριά}

\ex Γιατί ένας αριθμός που τελειώνει σε 0 διαιρείται και με το 2 και με το 5; (Βοήθεια: $10=2\cdot 5$.)

\ex Γράψε το 1.452 ως $1\cdot 999+4\cdot 99+5\cdot 9+(1+4+5+2)$ και εξήγησε με τη βοήθεια αυτής της γραφής αν διαιρείται με το 3 και αν διαιρείται με το 9.

\ex Ποιος είναι ο μικρότερος τετραψήφιος αριθμός που διαιρείται ταυτόχρονα με τα 2, 3, 4, 5 και 9;

\ex Οι αριθμοί 121, 132, 143 και 154 διαιρούνται όλοι με το 11. Τι παρατηρείς στα ψηφία τους; Δοκίμασε το ίδιο μοτίβο στο 165 και στο 175 και έλεγξε με διαίρεση.

\ex Ο Στέλιος ισχυρίζεται ότι ο 2.026 διαιρείται με το 3, γιατί τελειώνει σε 6, και το 6 διαιρείται με το 3. Ποιο είναι το λάθος του; Τι θα έπρεπε να είχε κάνει;
